% ===== tr-07 块2：第 5 章 pp.5--7（原文页 5--7）=====

% [原文 p.5]
\begin{itemize}
\item 一族 Einstein 度量的行为强烈依赖于测地球体积是否满足非坍缩条件：存在一致的常数 $\kappa>0$ 使得
\[
\mathrm{Vol}_g(B_g(r))\ge\kappa r^{\dim_{\R}X},\qquad \forall\, 0<r\le\mathrm{diam}(X).
\]
在非坍缩情形下有良好的收敛与正则性理论 [13]。另一方面，坍缩情形下的度量退化在黎曼几何中大体上仍属未知领域；对一族坍缩的 Calabi--Yau 度量而言，半平坦度量渐近是一项高度不平凡的涌现特征。
\item Kähler 几何中反复出现的主题是度量与复几何之间的相互作用。非坍缩情形下的 Kähler--Einstein 度量与射影几何相联系 [23]。大复结构极限是一类非常剧烈的极化退化，其超越性质（与指数和对数相关）不能被传统的射影几何充分刻画；相反，非阿基米德几何凸显为自然的框架。于是人们可以问 Calabi--Yau 度量与非阿基米德几何之间的关系，而这个问题最终与 SYZ 猜想相关。
\end{itemize}

% [Tutor Note: 非坍缩条件中的指数 $\dim_{\R}X$ 是实维数（对 $n$ 维 CY 流形即 $2n$），$\kappa$ 要求在一族度量中一致；这条条件区分了“坍缩”与“非坍缩”，Gromov--Hausdorff 紧性理论只在非坍缩情形下给出良好的极限空间。]

\subsection{带椭圆曲面的坍缩 K3 曲面}

与 SYZ 猜想相关的一项有影响的早期工作是 Gross--Wilson [32] 关于带椭圆纤维化 $\pi:X\to\C\PP^1$ 的 K3 曲面 $X$ 上超 Kähler 度量的粘合描述。

由超 Kähler 旋转，特殊拉格朗日环面纤维可以看作另一个复结构下的全纯椭圆曲线。在一般情形下，该椭圆纤维化有 24 条 $I_1$ 型奇异纤维，也就是说，纤维化中的局部奇点在复几何上以 $(z_1,z_2)\mapsto z_1z_2$ 为模型。他们固定 K3 曲面上的 Kähler 类 $[\omega_X]$ 以及 $\C\PP^1$ 上的 $[\omega_{\PP^1}]$，并描述 Kähler 类 $\tau[\omega_X]+\pi^*[\omega_{\PP^1}]$（$0<\tau\ll1$）中的 Calabi--Yau 度量 $\omega_\tau$。若干关键的概念性特征是：

\begin{itemize}
\item 在一般区域中，度量 $\omega_\tau$ 在指数小的误差范围内以半平坦度量为模型，后者可以通过椭圆曲线的周期积分显式描述。半平坦度量渐近失效的 K3 曲面子集，其长度尺度为 $O(\tau^{1/2}|\log\tau|^{1/2})$。
\item 当 $\tau\to0$ 时，度量 $\omega_\tau$ 在 Gromov--Hausdorff 意义下收敛到 $\C\PP^1$ 上的一个奇异度量。K3 曲面的直径具有 $O(1)$ 的常数量级。一般椭圆曲线纤维的长度尺度为 $O(\tau^{1/2})$，当 $\tau\to0$ 时收缩到零尺寸。
\end{itemize}

% [Tutor Note: 所取的 Kähler 类 $\tau[\omega_X]+\pi^*[\omega_{\PP^1}]$ 中，交叉项给出 $\int_X(\cdot)^2$ 的 $O(\tau)$ 主部（$\int_X\omega_X\wedge\pi^*\omega_{\PP^1}=\int_{\PP^1}\left(\int_{X_b}\omega_X\right)\omega_{\PP^1}\ne0$），故 K3 的总体积随 $\tau\to0$ 趋于零，而直径保持 $O(1)$------这正是“坍缩”而非“缩成一点”的典型图像。]

% [Tutor Note: 24 条 $I_1$ 型奇异纤维与 K3 的 Euler 示性数 $\chi(K3)=24$ 相符：每条 $I_1$ 纤维（结点椭圆曲线）对 Euler 示性数贡献 $1$。]

% [Tutor Note: 原文 p.5 的失效长度尺度 $O(\tau^{1/2}|\log\tau|^{1/2})$ 与 p.6 的 Taub--NUT 区域尺度 $O(\tau^{1/2}|\log\tau|^{-1/2})$ 只差 $|\log\tau|$ 的指数符号（$+1/2$ 对 $-1/2$），极易抄混；已分别以 700 dpi 裁图核实（p.5 y 632--652、p.6 y 262--282 PDF 点）。]

% [原文 p.6]
\begin{itemize}
\item 在 $I_1$ 型奇异纤维的邻域内，度量以 Ooguri--Vafa 度量为模型，后者是借助 Gibbons--Hawking ansatz 构造的显式 $S^1$ 不变非完备超 Kähler 度量。
\item Ooguri--Vafa 度量的渐近几何与半平坦度量相配。这是粘合构造的基本要求。
\item 在 $I_1$ 纤维的结点奇点附近，Ooguri--Vafa 度量包含一个用小尺度 Taub--NUT 度量近似的区域，其长度尺度为 $O(\tau^{1/2}|\log\tau|^{-1/2})$。这些区域集中了 K3 曲面几乎全部的 $L^2$-黎曼曲率。
\end{itemize}

Gross--Wilson 图景代表了 SYZ 猜想上的\emph{最佳期望}，\footnote{文献中有时被忽略的一个告诫是：Gross--Wilson 情形下的超 Kähler 旋转并不完全是一族极化退化，因此并不完全符合我们的大复结构极限概念。在我们的视角下，Gross--Wilson 是一种启发，而不是 SYZ 猜想的例子。}它包含对度量的近乎显式的描述，并且 $X$ 上整体存在特殊拉格朗日环面纤维化。这一图景被部分地推广到带有全纯拉格朗日阿贝尔簇纤维化的高维超 Kähler 流形。经过超 Kähler 旋转，这些可以看作特殊拉格朗日纤维化。Tosatti 等人 [72][31] 建立了一般区域中的半平坦度量渐近，以及到基流形的 Gromov--Hausdorff 坍缩。

\subsection{Calabi--Yau 三维流形上的最佳期望}

关于本节的背景资料，见 [42, 第 8、9 章] 以及 [54] 的引言。在 SYZ 猜想提出后的最初几年里，人们基于与超 Kähler 情形的类比、以及拓扑与复几何方面的考量，持有一种过于乐观的信念：

\begin{itemize}
\item 特殊拉格朗日纤维化整体存在，且由一个 $C^\infty$ 映射给出，尽管某些纤维可能是奇异的。
\item 基上的判别轨迹是余维二的。对 Calabi--Yau 三维流形，在适当的生成性假设下，判别轨迹是一个三价图，有两类顶点，称为正顶点与负顶点。\footnote{“正/负顶点”这两个名称来自环面纤维化的一些老式拓扑模型，其中最奇异的纤维分别具有 Euler 示性数 $\pm1$。}
\item 沿着三价图的边，SYZ 纤维化的奇点在横向上以 $I_1$ 奇点为模型。
% [原文 p.7]
\item 正顶点附近的局部区域在复几何上是
\[
\{z_0z_1z_2 = 1-z_3\}\subset\C^3\times\C^*_{z_3},\qquad \Omega\propto\frac{1}{z_3}dz_0\wedge dz_1\wedge dz_2
\]
中的一个大开子集；而负顶点区域则以
\[
\{z_3z_4 = 1-z_1-z_2\}\subset(\C^*)^2_{z_1,z_2}\times\C^2_{z_3,z_4},\qquad \Omega\propto\frac{1}{z_1z_2}dz_2\wedge dz_3\wedge dz_4
\]
中的一个大开子集为模型。
\end{itemize}

% [Tutor Note: 负顶点的体积形式写作 $\frac{1}{z_1z_2}dz_2\wedge dz_3\wedge dz_4$，分子从 $dz_2$ 而不是 $dz_1$ 开始，看似不对称；但该模型所在的图以 $z_2,z_3,z_4$ 为坐标（$z_1=1-z_2-z_3z_4$），故这是该图上的非零全纯 3-形式。600 dpi 裁图（p.7，y 145--260 PDF 点）确认原文确为 $dz_2\wedge dz_3\wedge dz_4$，并非笔误。]

这种天真设想受到了 Joyce [41] 的挑战，依据是他关于特殊拉格朗日奇点的观察------这些奇点与全纯纤维化中所见的奇点有显著不同。紧 Calabi--Yau 流形在大复结构附近若真的整体存在特殊拉格朗日纤维化，人们预期它具有微妙得多的结构：

\begin{itemize}
\item 特殊拉格朗日纤维化通常不是由 $C^\infty$ 映射给出的，而最多是分段光滑的。
\item 基上的判别轨迹通常不是余维二的；相反，人们预期三价图被加厚为余维一的“带状”。加厚的量在大复结构极限下很可能趋于零。
\item $I_1\times\R$ 局部奇点模型不能为奇异的特殊拉格朗日纤维导出一个 Fredholm 形变理论，因此应当由其他奇点模型取代。
\end{itemize}

% [Tutor Note: $I_1\times\R$ 模型即“横截方向为 $I_1$（结点椭圆曲线）奇点、沿纤维方向为 $\R$”的局部模型。原文指出它给不出 Fredholm 形变理论，因而必须更换奇点模型------这是 Joyce 的观察所带来的结构性困难，而不是技术性细节。]

撇开特殊拉格朗日局部奇点这一实质性困难不谈，另一个主要困难是在复维数三的情形下理解大复结构极限附近的 Calabi--Yau 度量。乐观的期望受到复维数二情形下 Gross--Wilson 图景的启发。按假设，Calabi--Yau 三维流形在 Gromov--Hausdorff 意义下接近于一个实三维流形，人们相信其拓扑是三维球面，\footnote{这一期望来自本质骨架的拓扑，见下面第 5.3.3 节。}其中含有三价图，使得

\begin{itemize}
\item 远离三价图处，Calabi--Yau 度量在指数小的误差范围内是半平坦的。
\item 横向于三价图中的边，度量以 Gross 与 Wilson 图景中出现的 Ooguri--Vafa 度量为模型。
\item 在三价图的正、负顶点附近，局部度量以 Ooguri--Vafa 度量的某种推广为模型。
\end{itemize}
